Nel lavoro di tesi vogliamo gestire una competizione CTF (Capture The
Flag).

Di solito in una competizione CTF, uno o piu' utenti (eventualmente
organizzati in team) si sfidano per risolvere un'insieme di challenge,
ossia recuperare una \textbf{flag}: una stringa del tipo
\passthrough{\lstinline!\{hai\_trovato\_la\_flag\}!}.

Ci concentreremo nel lavoro in un tipo di competizione CTF, basato sulla
gestione di uno o piu sistemi vulnerabili da mitigare, dunque una
competizione di Difesa dove la flag si ottiene unicamente difendendo il
sistema, ossia \emph{``patchando''} la vulnerabilita'. Nel lavoro si
tiene conto del fatto che il sistema vulnerabile spesso potrebbe essere
costituito da uno o piu host/nodi che comunicano tra di loro.

Le CTF hanno lo scopo di migliorare la formazione degli specialisti del
settore, tramite la simulazione di scenari realistici, complessi o
quantomeno molto specifici. In particolare e' fondamentale formare
studenti e le risorse aziendali, per la gestione e mitigazione di
sistemi vulnerabili, in modo da prevenire perdite dovute ad attacchi
informatici in scenari reali. La simulazione di questi scenari avviene
all'interno di un ambiente denominato Cyber Range.

Abbiamo dunque sviluppato \emph{BlueAgent}, un'orchestratore di sistemi
vulnerabili basato sulla tecnologia dei container che implementa il
Cyber Range di cui abbiamo bisogno per gestire competizioni di Difesa.
BlueAgent si occupa inoltre di verificare che le vulnerabilita presenti
sui container siano state corrette, tramite un sistema di verifica che
avviene completamente all'oscuro dei giocatori, ossia la colonazione dei
container in un'ambiente separato.

INSERIRE RISULTATI OTTENUTI e METODOLOGIA USATA
